\noindent\textbf{Part 1 (15 points): Einstein Ring}

\begin{description}
\item[(a)] The redshift $z$ measures the shift in the wavelength of radiation
  emitted by a source ($\lambda$) and observed at a longer wavelength by the
  observer ($\lambda + \Delta\lambda$); it is caused by the expansion of the
  Universe between the source and the observer. For nearby objects, we have
  \[ z = \frac{\Delta\lambda}{\lambda} \simeq \frac{\Delta V}{c} \]
  \hfill(0.5 points)

  Using Hubble's law gives
  \[ \Delta V = H_0 d \]
  \hfill(0.5 points)

  where $d$ is the distance between the observer and the source. Hence,
  $d = zc/H_0$. This result leads to
  \[ d_S = \frac{0.406c}{H_0} = 1800\,\mathrm{Mpc} \]
  \hfill(0.5 points)
  \[ d_L = \frac{0.042c}{H_0} = 190\,\mathrm{Mpc} \]
  \hfill(0.5 points)

  \textit{Note: If students then calculate the angular diameter distance as
  $d/(1+z)$ they should also receive full points.}

\item[(b)] From Figure T09.2, and using the small angle approximation, we
  obtain
  \[ d_{LS} = d_S - d_L \]
  \[ \theta d_S = 2\alpha d_{LS} \]
  \hfill(1 point)
  \[ \theta d_L = b \]
  \hfill(1 point)

  Based on the relation $2\alpha = 2R^*/b$, we find
  \begin{align*}
    2\alpha &= \frac{2R^*}{\theta d_L} \\
    \theta d_S &= \frac{2R^* d_{LS}}{\theta d_L} \\
    \theta^2 &= \frac{2R^* d_{LS}}{d_L d_S} \\
    \theta &= \sqrt{\frac{2R^* d_{LS}}{d_L d_S}}
            = \sqrt{\frac{2R^*(d_S - d_L)}{d_L d_S}}
  \end{align*}
  \hfill(1 point)

  \textit{Note: If students include the minus sign in their answer, deduct
  0.5 points.}

\item[(c)] \textbf{Determine the distance of closest approach}

  The distance of closest approach is given by
  \[ b = \theta d_L \]
  Substituting $\theta = 2.5$ [arcsec] and $d_L = 190$ [Mpc], we get
  \[ b = 2.3\,\mathrm{kpc} \]
  \hfill(1 point)

  \textbf{Determine the mass of the lens galaxy}

  The Schwarzschild radius of an object of mass $M$ is
  \[ R^* = \frac{2GM}{c^2} \]
  \hfill(0.5 points)

  Using the expression derived in the previous part,
  \[ \theta^2 = \frac{2R^* d_{LS}}{d_L d_S} \]
  we obtain
  \[ R^* = \frac{\theta^2 d_L d_S}{2 d_{LS}} \]
  Hence, the mass of the lens galaxy is
  \[ M = \frac{R^* c^2}{2G} = \frac{\theta^2 d_L d_S}{2 d_{LS}}\,\frac{c^2}{2G} \]
  \hfill(0.5 points)

  Substituting $\theta = 2.5$ [arcsec] $= 12.1 \times 10^{-6}$ [rad],
  $d_L = 190$ [Mpc], and $d_S = 1800$ [Mpc], we get
  \[ M = 1.6 \times 10^{11}\,\mathrm{M}_\odot \]
  \hfill(1 point)

  \textit{Note that this is only the mass enclosed within the Einstein ring
  and is therefore smaller than the total mass of the lens galaxy.}

\item[(d)] When the lens is displaced by a small distance $\varepsilon$
  (Figure T09.3), the distance of closest approach decreases for the upper
  ray, increasing its deflection. The ray that reaches the observer must
  therefore be emitted at a larger angle. Conversely, for the lower ray, the
  distance of closest approach increases, reducing the deflection, so the
  observed ray must be emitted at a smaller angle.

  Let $\theta'$ be the new angular radius of the ring, namely the angle
  between the detected light ray and the line joining the observer and the
  lens, $2\alpha'$ be the new deflection angle, and $b'$ be the new distance
  to the closest approach.

  First, considering the upper part of the figure, we have
  \[ \theta + \frac{x}{d_L} = \theta' + \frac{\varepsilon}{d_L} \]
  \[ \Rightarrow x = \theta' d_L - \theta d_L + \varepsilon \qquad (1) \]
  Since $b = \theta d_L$ and $b' = \theta' d_L$, we obtain
  \[ x = b' - b + \varepsilon \qquad (2) \]
  \hfill(1 point)

  Repeating the calculation described in b), we find
  \[ 2\alpha' d_{LS} = \left(\theta' + \frac{\varepsilon}{d_L}\right) d_S
     \qquad (3) \]
  \hfill(0.5 points)

  Moreover,
  \[ 2\alpha' = \frac{2R^*}{\theta' d_L} \qquad (4) \]
  \hfill(0.5 points)

  Hence,
  \[ \left(\theta' + \frac{\varepsilon}{d_L}\right) d_S
     = \frac{2R^*}{\theta' d_L}\, d_{LS} \qquad (5) \]
  For the perfectly aligned configuration, we have
  \[ \theta d_S = \frac{2R^*}{\theta d_L}\, d_{LS} \qquad (6) \]
  Dividing (5) by (6) gives
  \[ \frac{\theta'}{\theta} + \frac{\varepsilon}{\theta d_L}
     = \frac{\theta}{\theta'} \]
  \[ \Rightarrow \frac{\theta'}{\theta} + \frac{\varepsilon}{b}
     = \frac{\theta}{\theta'}
     \ \text{ or }\
     \frac{b'}{b} + \frac{\varepsilon}{b} = \frac{b}{b'} \qquad (7) \]
  \hfill(1 point)

  To simplify the calculation, we define
  \[ \frac{\theta'}{\theta} = 1 + k\frac{\varepsilon}{b} \qquad (8) \]
  Since $\varepsilon/b \ll 1$, Eq.\ (7) reduces to
  \[ 1 + k\frac{\varepsilon}{b} + \frac{\varepsilon}{b}
     = 1 - k\frac{\varepsilon}{b}
     \Rightarrow k = -\frac{1}{2} \qquad (9) \]
  Substituting $k$ into Eq.\ (8), we obtain
  \[ \frac{\theta'}{\theta} = 1 - \frac{\varepsilon}{2b} \qquad (10) \]
  Hence,
  \[ \frac{b'}{b} = 1 - \frac{\varepsilon}{2b}
     \Rightarrow b' - b = -\frac{\varepsilon}{2} \qquad (11) \]

  \textit{Note: Alternatively, Eq.\ (7) can be solved as follows:}

  \textit{Let $t = \theta'/\theta = b'/b$, we have}
  \[ t^2 + \frac{\varepsilon}{b}\,t - 1 = 0 \]
  \textit{This equation has two solutions, but we retain only the positive
  one, which is physically meaningful. The positive solution is}
  \[ t = \frac{-\varepsilon/b + \sqrt{(\varepsilon/b)^2 + 4}}{2} \]
  \textit{Since $\varepsilon/b \ll 1$, we can keep only the first-order term
  in $\varepsilon/b$,}
  \[ t = 1 - \frac{\varepsilon}{2b}
     \Rightarrow b' - b = -\frac{\varepsilon}{2} \]

  Substituting this result into Eq.\ (2), we finally get
  \[ x = \frac{\varepsilon}{2} \qquad (12) \]
  \hfill(2 points)

  For the lower part of the figure, we obtain a similar result by changing
  $\varepsilon$ to $-\varepsilon$ in corresponding equations:
  \[ x' = \frac{\varepsilon}{2} \qquad (13) \]
  \hfill(1 point)

  The angular displacement of the ring image is
  \[ \Delta\theta = \frac{\varepsilon}{2 d_L} \qquad (14) \]
  \hfill(1 point)

  To first order, the light rays in the plane perpendicular to the plane of
  Figure T09.3 are unaffected. Hence, the observer sees the Einstein ring
  displaced in the same direction as the lens, but by only \textit{half the
  lens displacement}.
\end{description}

\medskip\noindent\textbf{Part 2 (5 points): Gravitational Delay}

\begin{description}
\item[(a)] Since the light ray approaching the Sun tangentially, the distance
  of closest approach is
  \[ b = R_\odot \]
  \hfill(0.5 points)

  Hence, the total deflection angle is
  \[ 2\alpha = \frac{2R^*}{b} = \frac{2R^*}{R_\odot} \]
  \hfill(0.5 points)

  The Schwarzschild radius of the Sun is given by
  \[ R^* = \frac{2G\mathrm{M}_\odot}{c^2} = 2.95\,\mathrm{km}\ (\sim 3\,\mathrm{km}) \]
  \hfill(0.5 points)

  Substituting the values of $R^*$ and $R_\odot$, we obtain
  \[ 2\alpha = \frac{2 \times 2.95}{6.96 \times 10^5}
     = 0.85 \times 10^{-5}\,[\mathrm{rad}] = 1.75\,[\mathrm{arcsec}] \]
  \hfill(0.5 points)

  \textit{(1.75--1.78 arcsec for full mark)}

\item[(b)] The relation $\mathrm{d}\delta/\mathrm{d}R = 2\alpha/c$ was
  derived using the small-angle geometry for light rays passing close to the
  Sun. This approximation becomes less accurate as the distance from the Sun
  increases. Nevertheless, for this simplified exercise, we use it to estimate
  the total gravitational delay.

  The gravitational delay experienced by a light ray between the Sun surface
  ($R = R_\odot$) and the Earth ($R = d_\odot = 1\,\mathrm{au}$) is the
  integral of the derivative between these limits, namely
  \[ \delta = \int_{R_\odot}^{d_\odot} \frac{2\alpha}{c}\,\mathrm{d}R \]
  \hfill(1 point)

  The total bending angle of a light ray grazing the Sun is
  \[ 2\alpha = \frac{2R^*}{R} \]
  Hence, the gravitational delay is
  \[ \delta = \int_{R_\odot}^{d_\odot} \frac{2R^*}{cR}\,\mathrm{d}R
     = \frac{2R^*}{c}\int_{R_\odot}^{d_\odot} \frac{1}{R}\,\mathrm{d}R \]
  \hfill(0.5 points)
  \[ \Rightarrow \delta = \frac{2R^*}{c}\ln\left(\frac{d_\odot}{R_\odot}\right) \]
  \hfill(1 point)
  \[ \Rightarrow \delta = \frac{2 \times 2.95}{3 \times 10^5}
     \ln\left(\frac{1.496 \times 10^8}{6.96 \times 10^5}\right)
     = 106\,\mu\mathrm{s} \]
  \hfill(0.5 points)

  \textit{Note: If students integrate from the solar radius $R_\odot$ to any
  $R' < 1\,\mathrm{au}$, they will also receive full credit.}
\end{description}
